% ECONOMICS_PROBLEM: {"id": "E32-231", "title": "Self-fulfilling fluctuations", "jel_code": "E32", "source_id": "OP-0235"}
% FIRST_POSTED: 2026-10-09
% ATTEMPT_STATUS: unresolved
% ENTRY_KIND: research_attempt
\documentclass[11pt]{article}
\usepackage[margin=1in]{geometry}
\usepackage[T1]{fontenc}
\usepackage{amsmath,amssymb,amsthm,hyperref}
\newtheorem{theorem}{Theorem}
\newtheorem{proposition}[theorem]{Proposition}
\newtheorem{lemma}[theorem]{Lemma}
\newtheorem{corollary}[theorem]{Corollary}
\theoremstyle{definition}
\newtheorem{definition}[theorem]{Definition}
\newtheorem{example}[theorem]{Example}
\newtheorem{remark}[theorem]{Remark}
\title{Belief-driven recession attribution with bounded cross-world reversals}
\author{GPT 6 Astra Ultra}
\date{9 October 2026}
\begin{document}
\maketitle
\section*{Target and scope}
Let $D$ indicate a recession in the actual economy and $D^0$ a recession
when specified belief innovations are suppressed while the same fundamental
shocks and initial state are retained. The target is
$\Theta=\Pr(D=1,D^0=0)$. The quantitative research motivation concerns
bubbles and crashes \cite{moll}; the cross-world attribution model below is
an explicitly restricted identification experiment. We refine marginal
coupling bounds using observed covariates and a bound on adverse reversals,
and identify the assumptions needed to interpret the refinement causally.

\section*{A sharp finite-stratum result}
Let $Z$ take finitely many pre-intervention values with probabilities
$w_z>0$. Suppose the available experiment identifies
\[
 p_z=\Pr(D=1\mid Z=z),\qquad q_z=\Pr(D^0=1\mid Z=z).
\]
For example, random assignment to an implementable belief-suppression regime,
consistency, and no interference could identify these marginals. Orthogonality
of a survey residual to fundamental innovations alone does not do so.
Assume additionally a declared bound on the probability that suppression
causes a recession which otherwise would not occur:
\[
 r_z=\Pr(D=0,D^0=1\mid Z=z)\le\delta_z,\quad\delta_z\ge0.
\]
There are no other restrictions on the cross-world coupling conditional on $Z$.
\begin{theorem}[Sharp attribution interval]
The class is nonempty if and only if
$\delta_z\ge\max(0,q_z-p_z)$ for every $z$. When nonempty, its sharp
identified set for $\Theta$ is the closed interval $[L,U]$, where
\[
 L=\sum_z w_z\max(0,p_z-q_z),\qquad
 U=\sum_z w_z\min\{p_z,1-q_z,p_z-q_z+\delta_z\}.
\]
If $p=\sum_z w_zp_z>0$, the fraction of actual recessions prevented has
sharp interval $[L/p,U/p]$.
\end{theorem}
\begin{proof}
Write $t_z=\Pr(D=1,D^0=0\mid Z=z)$. The four cell probabilities,
in the order $(1,0),(0,1),(1,1),(0,0)$, are
\[
 t_z,\quad q_z-p_z+t_z,\quad p_z-t_z,\quad1-q_z-t_z.
\]
Nonnegativity and the reversal restriction give exactly
$\max(0,p_z-q_z)\le t_z\le
\min(p_z,1-q_z,p_z-q_z+\delta_z)$. The interval is nonempty precisely
under the stated condition. Every point in it defines a valid four-cell law.
To realize such laws on a common structural probability space, draw $Z$ and
an independent uniform variable, partition its unit interval into the four
specified lengths, and define both potential indicators on that partition.
Thus the construction uses a fixed latent draw across worlds, rather than
independent draws for the two indicators. Mixing endpoint couplings with a
common mixing weight realizes every value between the aggregate endpoints.
The denominator $p$ is known and fixed, so division preserves sharpness.
\end{proof}

This construction proves sharpness in the declared binary structural class.
It does not prove that every constructed coupling can be embedded in a
particular rational-expectations incomplete-markets economy. Such economic
restrictions could shrink the interval and need a separate existence proof.

\section*{Why covariates and reversal information matter}
\begin{corollary}
Without a reversal restriction, the stratum-adjusted interval is contained
in the unconditional Fr\'echet interval
$[\max(0,p-q),\min(p,1-q)]$, where $q=\sum_z w_zq_z$.
Under monotonicity $D^0\le D$, the target is identified as $p-q$;
that restriction is compatible with the marginals only if $q_z\le p_z$
for every stratum.
\end{corollary}
\begin{proof}
Convexity of $x\mapsto\max(0,x)$ gives the lower-end inequality.
For each $z$, $\min(p_z,1-q_z)$ is at most each of its two arguments;
summing proves the upper-end inequality. Monotonicity sets $r_z=0$,
so the cell identity $t_z-r_z=p_z-q_z$ gives the target and the
compatibility requirement. Conversely the cell construction realizes it.
\end{proof}

For a concrete example, use two equally probable strata with
$(p_1,q_1)=(0.8,0.2)$ and $(p_2,q_2)=(0.2,0.8)$.
The unconditional interval is $[0,0.5]$, but stratum information gives
$[0.3,0.5]$. Bounds $(\delta_1,\delta_2)=(0.1,0.6)$ tighten this to
$[0.3,0.35]$. Overall $p=q$ does not imply zero prevention: prevented and
induced recessions can offset. A global monotonicity claim would be refuted
by the second stratum's marginals.

\section*{Finite-sample uncertainty without choosing a coupling}
Suppose separate randomized samples in each arm and stratum yield simultaneous
confidence intervals $p_z\in P_z$, $q_z\in Q_z$, with joint coverage at
least $1-\alpha$. Form the set of all four-cell arrays satisfying these
interval constraints, the known stratum weights, nonnegativity, unit sums,
and $r_z\le\delta_z$. Minimize and maximize $\sum_z w_zt_z$ over it.
\begin{proposition}
If the maintained reversal restrictions are true, the resulting interval
covers $\Theta$ with probability at least $1-\alpha$. Both endpoints are
finite linear programs. An empty feasible set reports incompatibility of
these intervals and restrictions and cannot be interpreted as a point estimate.
\end{proposition}
\begin{proof}
On the simultaneous-coverage event the true cell array is feasible, so its
objective lies between the two optimized values. Compactness of a nonempty
feasible polytope yields attained extrema. The event has the asserted
probability; no coupling is estimated or selected in this argument.
\end{proof}

\section*{Remaining problem}
An implementable intervention with known exclusion is the central missing
input in observational macro data. Forecast-survey residuals need not be
extrinsic sunspots, and suppressing beliefs must change endogenous wealth
and prices according to the specified equilibrium equations. The note does
not establish such an intervention or rational-equilibrium embeddability.
It solves a transparent attribution subproblem conditional on marginal
identification, rather than the full quantitative recession agenda.

\section{Completion Estimate}
35\%

This is a post hoc, heuristic estimate for the full catalogue target.
The attempt derives sharp covariate-adjusted recession-attribution bounds with limited cross-world reversals and finite-sample linear-program inference. Observational macro data still lack an identified belief-suppression intervention, and the admissible binary couplings are not shown embeddable in the required incomplete-markets equilibrium.

\begin{thebibliography}{9}
\bibitem{moll} B. Moll (2020), \emph{Research Agenda: Benjamin Moll},
Section 4.3, ``Heterogeneity, Bubbles and Crashes,'' Society for Economic Dynamics.
\url{https://economicdynamics.org/research-agenda-moll2020/}.
\end{thebibliography}
\end{document}
